% =====================================================================
%  Script de soutenance (15 min) — MCP-Powered BI Ecosystem
%  Partie 1 : Mohamed (slides 1 à 8) · Partie 2 : Zakariae (slides 9 à 15)
%  Compatible pdfLaTeX (Overleaf) ET XeLaTeX / Tectonic.
% =====================================================================
\documentclass[11pt,a4paper]{article}

\usepackage{iftex}
\ifPDFTeX
  \usepackage[utf8]{inputenc}
  \usepackage[T1]{fontenc}
  \usepackage{lmodern}
\else
  \usepackage{fontspec}
\fi
\usepackage[french]{babel}
\usepackage[a4paper,margin=2.2cm,headheight=14pt]{geometry}
\usepackage{setspace}
\onehalfspacing
\usepackage{microtype}
\usepackage{enumitem}
\setlist{topsep=2pt,itemsep=2pt,leftmargin=1.6em}
\usepackage{needspace}

% ---- Couleurs (charte teal du projet) ------------------------------
\usepackage{xcolor}
\definecolor{accent}{HTML}{0D9488}
\definecolor{accentdark}{HTML}{047857}
\definecolor{accentsoft}{HTML}{D1FAF2}
\definecolor{ink}{HTML}{0F2A30}
\definecolor{muted}{HTML}{5B7178}

\usepackage{fancyhdr}
\pagestyle{fancy}\fancyhf{}
\fancyhead[L]{\small\color{muted}Script de soutenance}
\fancyhead[R]{\small\color{muted}MCP-Powered BI Ecosystem}
\fancyfoot[C]{\small\color{muted}\thepage}
\renewcommand{\headrulewidth}{0.4pt}

\usepackage[hidelinks]{hyperref}

% ---- Bandeau « partie / présentateur » -----------------------------
\newcommand{\presheader}[3]{%
  \bigskip\par\noindent
  \colorbox{accentdark}{%
    \parbox[c][0.98cm][c]{\dimexpr\linewidth-2\fboxsep\relax}{%
      \hspace{4pt}{\large\bfseries\color{white}#1}\quad
      {\color{accentsoft}\normalsize #2}\hfill
      {\color{accentsoft}\small #3}\hspace{4pt}}}%
  \par\medskip}

% ---- En-tête d'une slide -------------------------------------------
\newcommand{\slidecue}[3]{%
  \needspace{4\baselineskip}%
  \medskip\par\noindent
  {\color{accentdark}\rule[-3pt]{4pt}{17pt}}\hspace{8pt}%
  {\large\bfseries\color{ink}Slide #1\, — \,#2}%
  \hfill{\small\color{muted}#3}\par
  \nopagebreak\smallskip}

% ---- Didascalie (indication scénique) ------------------------------
\newcommand{\cue}[1]{{\itshape\color{muted}#1}}

\begin{document}
\thispagestyle{fancy}

% ---- Titre ----------------------------------------------------------
\begin{center}
  {\Large\bfseries\color{ink}Script de présentation — Soutenance PFA}\\[3pt]
  {\large\color{accentdark}MCP-Powered BI Ecosystem}\\[8pt]
  {\color{muted}Durée cible : \textbf{15 minutes} \;\textbullet\;
   Mohamed Rouane \emph{(slides 1 à 8)} \;\textbullet\;
   Zakariae Trachni \emph{(slides 9 à 15)}}
\end{center}
\vspace{2pt}
\noindent\colorbox{accentsoft}{%
  \parbox{\dimexpr\linewidth-2\fboxsep\relax}{\small\color{ink}\hspace{2pt}
  \textbf{Conseils :} parlez lentement, regardez le jury, et changez de slide
  quand le texte correspondant est terminé. Les temps indiqués sont des repères
  (total $\approx$ 14 min, pour garder une marge). Le texte \cue{en italique gris}
  est une indication d'action, à ne pas lire à voix haute.\hspace{2pt}}}

% =====================================================================
\presheader{Partie 1}{Mohamed Rouane — slides 1 à 8}{$\approx$ 6 min 40}

\slidecue{1}{Page de garde}{$\approx$ 30 s}
Bonjour à toutes et à tous. Nous sommes Mohamed Rouane et Zakariae Trachni,
étudiants en Ingénierie Data Science et Cloud Computing à l'ENSA d'Oujda. Nous
vous présentons notre projet de fin d'année, encadré par le professeur
Abdelmounaim Kerkri. Son titre : \textbf{MCP-Powered BI Ecosystem}. C'est une
plateforme qui automatise l'analyse de données grâce à plusieurs agents
d'intelligence artificielle. Je commence par le contexte et la conception, puis
Zakariae présentera la réalisation et une démonstration en direct.

\slidecue{2}{Plan de la présentation}{$\approx$ 20 s}
Voici notre plan. Nous posons d'abord le contexte et le problème, puis nos
objectifs et les technologies utilisées. Ensuite, je détaille la conception :
l'architecture, les sept agents et le moteur d'analyse. Zakariae enchaînera
avec la réalisation, une démonstration en direct, l'étude de cas et les
résultats, et enfin la conclusion.

\slidecue{3}{Contexte \& problématique}{$\approx$ 1 min 10}
Commençons par le contexte. Pour transformer des données en décisions, on suit
toujours la même chaîne : on récupère les données brutes, on les valide, on les
nettoie, on calcule des indicateurs, puis on produit un rapport pour décider.
Le problème, c'est que cette chaîne est aujourd'hui presque entièrement
\textbf{manuelle}. Elle dépend fortement d'un analyste, qui devient un goulot
d'étranglement. C'est un travail répétitif, à refaire à chaque nouveau fichier.
Les solutions existantes sont souvent spécialisées pour un seul cas d'usage,
avec des hypothèses figées. Et la rédaction du rapport reste entièrement
humaine. D'où notre \textbf{problématique} : comment analyser automatiquement,
et de manière fiable, n'importe quel jeu de données, en combinant la rigueur
d'un programme et le discernement d'agents d'intelligence artificielle ?

\slidecue{4}{Objectifs du projet}{$\approx$ 40 s}
Nos objectifs découlent de ce problème. Premièrement, construire une
architecture à plusieurs agents, coordonnés par un serveur MCP. Deuxièmement,
un moteur capable de s'adapter à n'importe quelles données, sans rien supposer
du métier. Troisièmement, produire automatiquement deux livrables : un tableau
de bord interactif et un rapport décisionnel. Et quatrièmement, garantir la
fiabilité des résultats. Le point le plus important est en bas de l'écran :
\textbf{zéro hypothèse métier}. Le système fonctionne sur n'importe quel fichier
CSV, Excel ou JSON.

\slidecue{5}{Les briques technologiques}{$\approx$ 1 min}
Sur quelles technologies nous appuyons-nous ? D'abord les grands modèles de
langage, les LLM : ils savent raisonner et résumer, mais ils ont un défaut
connu — ils inventent parfois des informations. C'est ce qu'on appelle les
\og hallucinations \fg{}. Ensuite, les agents d'IA : un agent observe, décide,
puis agit, et plusieurs agents spécialisés coopèrent. Le troisième élément est
le c\oe ur de notre projet : le \textbf{protocole MCP}, un standard ouvert qui
fournit aux modèles un registre d'outils avec un contrôle d'accès. Enfin, côté
technique, nous utilisons Python, pandas, Plotly et FastAPI.

\slidecue{6}{Architecture en couches}{$\approx$ 1 min}
Passons à la conception. Notre système est organisé en couches, du haut vers le
bas. En haut, la couche présentation : l'interface web où l'utilisateur dépose
son fichier et consulte les résultats. En dessous, l'orchestration, qui pilote
tout le processus. Puis la couche des sept agents. Au centre, le \textbf{serveur
MCP}, qui contrôle tous les accès aux outils. En dessous, le moteur
déterministe, qui fait les calculs. Et tout en bas, le magasin d'artefacts, où
sont stockés le tableau de bord et le rapport. L'avantage : chaque couche est
indépendante, et le serveur MCP est le seul point de passage vers les outils.

\slidecue{7}{Sept agents, un serveur MCP}{$\approx$ 1 min}
Voici nos sept agents : l'Orchestrateur, l'Ingestion, la préparation des
données, le KPI, le Dashboard, la Publication et le Reporter. Chacun a un rôle
précis, et tous passent par le serveur MCP, qui expose quatorze outils. Le
principe clé est celui du \textbf{moindre privilège} : chaque agent ne peut
utiliser que les outils strictement nécessaires à son rôle ; toute demande hors
de son périmètre est refusée par défaut. Par exemple, à droite, l'agent KPI ne
peut appeler que les outils liés aux indicateurs : profiler les données,
suggérer et calculer les KPIs, détecter les anomalies. C'est une vraie garantie
de sécurité et de robustesse.

\slidecue{8}{Le moteur déterministe adaptatif}{$\approx$ 1 min}
Voici le moteur qui réalise le vrai travail d'analyse, en quatre étapes.
Premièrement, le \textbf{profilage} : il devine le rôle de chaque colonne —
est-ce une mesure, une date, un pays, un identifiant ? Deuxièmement, il calcule
les indicateurs adaptés à ces rôles. Troisièmement, il génère les graphiques :
tendances, barres, carte, corrélations. Quatrièmement, il extrait des insights,
comme les anomalies et les comparaisons. Le message important est en bas : la
\textbf{bonne division du travail}. Le moteur garantit l'exactitude des
chiffres ; les agents IA apportent le jugement et la rédaction. Conséquence
directe : \textbf{aucun chiffre n'est inventé par le modèle}. Je laisse
maintenant la parole à Zakariae. \cue{(Transition vers Zakariae.)}

% =====================================================================
\presheader{Partie 2}{Zakariae Trachni — slides 9 à 15}{$\approx$ 7 min 30}

\slidecue{9}{Fiabilité : la confiance avant tout}{$\approx$ 45 s}
Merci Mohamed. Je vais parler de la fiabilité, qui est au c\oe ur de notre
projet. Nous avons mis en place trois mécanismes. D'abord, les \textbf{filets de
sécurité} : si un agent oublie une étape, le pipeline la réalise à sa place ;
les livrables sont donc toujours complets. Ensuite, la \textbf{vérification des
chiffres} : chaque nombre du rapport est comparé aux indicateurs réellement
calculés. C'est notre garde-fou contre les hallucinations. Enfin, le
\textbf{suivi des coûts} : on mesure les jetons, le temps et le coût de chaque
agent. Grâce à ces filets, une exécution réussit toujours.

\slidecue{10}{Démonstration en direct}{$\approx$ 3 min 30}
Passons maintenant à la pratique, avec une démonstration en direct. Je vais
faire quatre choses :
\begin{enumerate}[label=\textbf{\arabic*.}]
  \item Je \textbf{téléverse} un fichier de données — ici un simple CSV.
  \item On \textbf{suit le pipeline} en direct : vous voyez les sept agents
        travailler les uns après les autres.
  \item Une fois terminé, j'ouvre le \textbf{tableau de bord} interactif, avec
        ses graphiques et la carte mondiale.
  \item Enfin, j'ouvre le \textbf{rapport décisionnel}, dont tous les chiffres
        ont été vérifiés.
\end{enumerate}
\cue{(Prenez votre temps : commentez ce qui s'affiche à l'écran à chaque étape.
Si la démonstration en direct pose problème, les aperçus à droite de la slide
montrent ces mêmes étapes — servez-vous-en comme repli.)}

\slidecue{11}{Étude de cas : matchs internationaux}{$\approx$ 45 s}
Pour valider notre approche, prenons un vrai jeu de données : près de
\textbf{dix-huit mille matchs} de football internationaux, de 1872 à 2024,
quatre-vingt-dix tournois et plus de deux cents équipes. Point important : nous
n'avons donné \textbf{aucune configuration} ; le système a tout découvert seul.
Il a reconnu deux mesures — les buts à domicile et à l'extérieur. Il a compris
que les équipes étaient des pays, ce qui a permis de générer une carte mondiale.
Et il a repéré la colonne de dates ainsi que les différentes dimensions, comme
le tournoi ou le stade.

\slidecue{12}{Des insights pertinents, automatiquement}{$\approx$ 1 min}
Et les résultats sont parlants. À gauche, le système a détecté tout seul une
\textbf{anomalie} : le match Australie 31 à 0 face aux Samoa américaines, en
2001 — c'est un record du monde. Il l'a isolé sans aucune connaissance du
football. À droite, une \textbf{comparaison de segments} : les équipes marquent
en moyenne 1,78 but à domicile, contre 1,48 à l'extérieur. Autrement dit,
l'avantage du terrain, mesuré quantitativement. Et je le redis : tous ces
chiffres présents dans le rapport ont été \textbf{vérifiés automatiquement}.

\slidecue{13}{Validation : 22 tests automatisés}{$\approx$ 30 s}
Côté qualité, nous avons écrit \textbf{vingt-deux tests automatisés}, et ils
sont tous au vert. Ils couvrent le profilage des colonnes, les permissions du
serveur MCP, la vérification des chiffres, la relance des agents, la détection
géographique et les insights. Ces tests ne font pas appel aux modèles de
langage : ils sont donc rapides et reproductibles.

\slidecue{14}{Conclusion \& perspectives}{$\approx$ 45 s}
En conclusion, nous avons construit une plateforme de Business Intelligence
\textbf{généraliste et automatisée}, qui fonctionne de bout en bout. Elle
combine un moteur déterministe et sept agents coordonnés par MCP, produit un
tableau de bord et un rapport reproductibles, le tout avec de vrais mécanismes
de fiabilité. Et elle a été validée sur près de dix-huit mille matchs réels.
Comme perspectives, nous envisageons l'interrogation en langage naturel, des
connecteurs vers des bases de données, des analyses prédictives, et
l'industrialisation de la plateforme.

\slidecue{15}{Merci}{$\approx$ 15 s}
Voilà qui conclut notre présentation. Nous vous remercions pour votre attention,
et nous sommes prêts à répondre à vos questions. \cue{(Rester détendu et à
l'écoute pour les questions du jury.)}

\end{document}
